%% kapa: exam=1 round=32 year=2021 period=1 form=A subject=economics subject_ko=경제학원론 seq=3 number=43
%% source: Q-Net 감정평가사 기출문제 (cst00309) artlSeq=5208397 file=extracted/1008/5208397/A형/2021년도 제32회 감정평가사 1차 1교시 A형.pdf
%% conversion: kordoc 4.24.51 -> tools/gen1.py (auto); review=proofread-sonnet; flags=none
\begin{question}{43}
수요곡선에 관한 설명으로 옳지 않은 것은?
\begin{choices}
  \item 우하향하는 수요곡선의 경우, 수요의 법칙이 성립한다.
  \item 기펜재(Giffen goods)의 수요곡선은 대체효과보다 소득효과가 크기 때문에 우하향한다.
  \item 사적재화의 시장수요는 개별수요의 수평 합이다.
  \item 우하향하는 수요곡선의 높이는 한계편익이다.
  \item 소비자의 소득이 변하면 수요곡선이 이동한다.
\end{choices}
\end{question}
\answer{2}
